\begin{recipe}[preparationtime = {\unit[20]{Min}},bakingtime = {\unit[70]{Min}},bakingtemperature={\protect\bakingtemperature{fanoven=\unit[160]{\textcelsius}}},portion = 1 Springform]{Himbeertarte}\label{himbeertarte}
	\ingredients{
		\multicolumn{2}{l}{\textbf{Teig:}}\\
		\unit[135]{g} & Weizenmehl\index{Mehl!Weizen}\\
		\unit[100]{g} & zimmerwarme Butter\index{Butter}\\
		\unit[50]{g} & Rohrzucker\index{Zucker!Rohr}\\
		\unit[1]{} & Ei (Größe M)\index{Ei}\\
		\unit[1]{TL} & Vanillezucker{Butter}\index{Zucker!Vanille}\\
		\unit[1]{TL} & Zitronenschale ungespritzt\index{Zitrone!Schale}\\
		\unit[1]{Prise} & Salz\\
		\multicolumn{2}{l}{\textbf{Füllung \& Baiser}}\\
		\unit[2]{Blatt} & Gelatine\index{Gelatine}\\
		\unit[100]{g} & Kuvertüre (Vollmilch)\index{Kuvertüre}\\
		\unit[4]{} & Eier (Größe M)\index{Ei}\\
		\unit[425]{g} & Himbeeren\index{Himbeere}\\
		\unit[30]{g} & Rohrzucker\index{Zucker!Rohr}\\
		\unit[120]{g} & Puderzucker\index{Zucker!Puder}
	}
	\preparation{
		\step Teigzutaten verkneten, zu einer Kugel formen, in Frischhaltefolie wickeln und \unit[60]{Minuten} in den Kühlschrank stellen.
		\step  2/3 des Teiges auf dem Springformboden ausrollen und mit einer Gabel mehrmals einstechen. Springformrand darumlegen.
		\step Übrigen Teig zu einer Rolle formen. Teigrolle als Rand auf den Teigboden legen und so an die Form drücken, dass ein \unit[3]{cm} hoher Rand entsteht. 
		\step \textbf{Eier trennen}. Eiweiß mit Puderzucker steif schlagen. 
		\step Himbeeren auf dem Boden verteilen. Jetzt weiß ich nicht was man hier machen soll...Normalerweise alles verrühren aber das geht mit der Gelatine nicht so...Im Zweifel an der Birnenwähe orientieren (s. Seite \pageref{birnenwaehe}).
		\step Die Masse auf dem Tarteboden verteilen und Häufchen aus dem Eischnee auflegen. Alles in den Ofen bei \unit[160]{\textcelsius} (Umluft) abbacken bis es schön braun ist.
	}
\end{recipe}
